\chapter{Kết luận và hướng phát triển}

\section{Kết quả đạt được}

Sau quá trình nghiên cứu, thiết kế và triển khai, nhóm đã xây dựng thành công hệ thống ChatBot hỗ trợ học tập ứng dụng mô hình Retrieval-Augmented Generation (RAG). Hệ thống cho phép người dùng đặt câu hỏi liên quan đến tài liệu môn học, đồng thời cung cấp môi trường Benchmark chuyên sâu để thực hiện đo lường, phân tích và so sánh chất lượng câu trả lời giữa các chiến lược phân mảnh dữ liệu (Chunking) và các mô hình biểu diễn từ (Embedding) khác nhau.

Các mục tiêu chính của đề tài đều đã được hoàn thành một cách trọn vẹn. Về mặt kiến trúc, nhóm đã phát triển thành công Backend bằng FastAPI đảm bảo khả năng xử lý bất đồng bộ mượt mà, đồng thời thiết kế giao diện người dùng bằng ReactJS trực quan và thân thiện. Về mặt xử lý dữ liệu AI, hệ thống đã tích hợp thành công cơ sở dữ liệu vector ChromaDB để phục vụ lưu trữ và truy xuất siêu tốc, hoàn thiện toàn bộ luồng xử lý RAG từ việc trích xuất file PDF cho đến sinh Prompt tối ưu. Hơn nữa, việc tích hợp thành công công cụ đo lường tự động (Benchmark) giúp hệ thống không chỉ là một ứng dụng hỏi đáp thông thường, mà còn là một công cụ nghiên cứu mạnh mẽ. 

Ngoài ra, toàn bộ quy trình làm việc nhóm đã được chuẩn hóa và quản lý hiệu quả thông qua việc ứng dụng phương pháp Agile Scrum trên Jira, cũng như việc kiểm soát phiên bản mã nguồn chặt chẽ bằng GitHub. Tài liệu đặc tả yêu cầu (SRS) và toàn bộ báo cáo này cũng được hoàn thiện chỉn chu bằng LaTeX.

\section{Ưu điểm}

Hệ thống đạt được nhiều ưu điểm nổi bật cả về mặt kỹ thuật lẫn ứng dụng thực tiễn:
\begin{itemize}
	\item \textbf{Tính chính xác cao và hạn chế Hallucination}: Nhờ ứng dụng triệt để cơ chế RAG, mọi câu trả lời của AI đều được neo chặt vào các đoạn tài liệu có thật do sinh viên hoặc giảng viên cung cấp, kèm theo trích dẫn nguồn rõ ràng.
	\item \textbf{Công cụ Benchmark đo lường định lượng}: Đây là điểm sáng của đề tài, giúp người dùng không phải đoán mò mà có thể dựa vào các chỉ số khách quan (Accuracy, Relevancy, Faithfulness) để chọn ra cấu hình RAG tốt nhất cho từng môn học.
	\item \textbf{Kiến trúc phân tầng và dễ mở rộng}: Việc phân tách rõ ràng Frontend, Backend và Database, kết hợp với thiết kế hướng đối tượng giúp hệ thống rất dễ bảo trì. Việc nâng cấp thêm các LLM mới (như chuyển từ Gemini sang GPT-4) có thể thực hiện nhanh chóng mà không làm xáo trộn kiến trúc.
	\item \textbf{Tốc độ truy xuất ấn tượng}: Sự kết hợp giữa FastAPI và ChromaDB mang lại thời gian phản hồi (Latency) cực ngắn cho các truy vấn phức tạp.
	\item \textbf{Quy trình phát triển chuyên nghiệp}: Toàn bộ quá trình code và thiết kế đều tuân thủ chặt chẽ vòng đời phát triển phần mềm (SDLC).
\end{itemize}

\section{Hạn chế}

Mặc dù đã đạt được những thành tựu đáng kể, hệ thống ở giai đoạn hiện tại vẫn tồn tại một số hạn chế nhất định cần được khắc phục:
\begin{itemize}
	\item \textbf{Giới hạn định dạng dữ liệu đầu vào}: Hiện tại hệ thống mới chỉ xử lý tốt các tài liệu định dạng văn bản chuẩn (như PDF, DOCX). Đối với các tài liệu chứa nhiều biểu đồ phức tạp, hình ảnh, hoặc công thức toán học chằng chịt, chất lượng trích xuất (OCR) đôi khi còn chưa đạt mức tối ưu nhất.
	\item \textbf{Môi trường triển khai}: Hệ thống hiện mới chỉ được chạy kiểm thử trên môi trường Local và server thử nghiệm quy mô nhỏ. Việc triển khai (Deployment) lên các nền tảng Cloud lớn (như AWS, Google Cloud) để phục vụ đồng thời hàng ngàn sinh viên chưa được thực hiện.
	\item \textbf{Chưa hỗ trợ quản lý định danh người dùng sâu (Authentication)}: Tính năng đăng nhập, phân quyền cá nhân hóa lịch sử hội thoại cho từng sinh viên cụ thể vẫn đang ở dạng thử nghiệm đơn giản.
\end{itemize}

\section{Hướng phát triển}

Từ những kết quả đã đạt được và các hạn chế còn tồn đọng, nhóm đề xuất các hướng phát triển tiếp theo nhằm hoàn thiện và nâng cấp hệ thống trong tương lai:
\begin{itemize}
	\item \textbf{Mở rộng hệ sinh thái mô hình AI}: Tích hợp và hỗ trợ đa dạng hơn các LLM mã nguồn mở (như Llama 3) chạy cục bộ, nhằm tiết kiệm chi phí gọi API và tăng tính bảo mật dữ liệu môn học.
	\item \textbf{Nâng cấp khả năng xử lý đa phương tiện (Multimodal)}: Ứng dụng các thuật toán xử lý ảnh tiên tiến để hệ thống có thể đọc hiểu và giải thích các biểu đồ, hình vẽ bên trong giáo trình.
	\item \textbf{Cá nhân hóa trải nghiệm học tập}: Tích hợp cơ chế đăng nhập (OAuth 2.0), cho phép theo dõi lịch sử và tiến độ học tập của từng sinh viên. Dựa trên dữ liệu hội thoại, hệ thống có thể chủ động đề xuất lộ trình ôn thi phù hợp.
	\item \textbf{Triển khai Cloud và Microservices}: Chuyển đổi kiến trúc sang dạng Microservices và đóng gói bằng Docker, triển khai lên nền tảng Kubernetes để đảm bảo hệ thống có khả năng tự động chịu tải (Auto-scaling) trong các kỳ thi cuối kỳ.
	\item \textbf{Phát triển ứng dụng Mobile}: Xây dựng thêm phiên bản ứng dụng di động (sử dụng React Native hoặc Flutter) giúp sinh viên thuận tiện hơn trong việc tra cứu kiến thức mọi lúc, mọi nơi.
\end{itemize}

\section{Tổng kết}

Đề tài đã hoàn thành xuất sắc việc xây dựng hệ thống ChatBot hỗ trợ học tập ứng dụng công nghệ Retrieval-Augmented Generation (RAG). Hệ thống không chỉ dừng lại ở mức độ một công cụ hỏi đáp tự động thông thường, mà còn cung cấp một nền tảng Benchmark toàn diện, cho phép đánh giá, phân tích và tối ưu hóa các chiến lược xử lý ngôn ngữ hiện đại. Những kết quả thu được từ đề tài này không chỉ giải quyết thiết thực nhu cầu tra cứu thông tin của sinh viên, mà còn tạo ra một bộ khung kiến trúc vững chắc, mở ra tiềm năng lớn cho việc tiếp tục nghiên cứu và phát triển các hệ thống trợ lý học tập thông minh ở quy mô lớn hơn trong tương lai.